\section{Путь А. Адаптивная кривизна и рабочий радиус}\label{sec:patha}

Путь А сохраняет рекурсию классического слайдинга. Он выбирает
кривизну модели по тому отрезку, который реально возник на пробном
шаге. При неудачной проверке меняется оценка конкретной пары.
Длины циклов и порядок пар пересчитываются перед следующей попыткой.

\subsection{Алгоритм}

Пусть $s=x_0$, $\delta=\rho/3$, $R_0>0$ и
$\lambda_i=\max\{\delta,\Lambda_i^0\}$ для начальных оценок
$\Lambda_i^0\ge0$. Начинаем с радиуса $R_0$ и кривизн
$\Lambda_i=\lambda_i$. На каждой попытке работаем на
$D=C\cap\ball_R(s)$. Сортируем пары по возрастанию
$h_i=\max\{1,\sqrt{\Lambda_i/\delta},M_i/\delta\}$. В формулах прохода
индексы относятся к этому порядку. Вводим
\begin{align}
 \alpha_0&=1,& \alpha_{t+1}&=\frac{2}{1+\sqrt{1+4/\alpha_t^2}},\\*
 T_1&=\ceil{2h_1},&
 T_i&=\ceil{16h_i/h_{i-1}}\qquad(i>1),\\*
 a_i&=\alpha_{T_i-1},&
 s_i&=\prod_{j=1}^{i}a_j,\qquad c_i=\Lambda_i+2M_i/s_i.
 \label{eq:pa-parameters}
\end{align}
Для списка моделей $m_j(z)=H_j\norm{z-u_j}^2/2+\ip{d_j}{z}$
лист возвращает $\Pi_D\bigl(\sum_j(H_j u_j-d_j)/\sum_j H_j\bigr)$.
Путь А использует евклидову метрику. В алгоритме~\ref{alg:path-a}
глубокие центры хранятся в $U$, как в алгоритме~\ref{alg:original}.
По сравнению с исходным проходом меняются расписание, кривизна
моделей и поправка центра: она использует $\Pi_D$. Проверки секущих
и радиуса могут вызвать повтор всего прохода.

\begin{algorithm}[htbp]
\caption{Путь А. Поиск кривизны и радиуса с проверяемым проходом}
\label{alg:path-a}
\begin{algorithmic}[1]
\REQUIRE $s\in C$, $\rho>0$, $R_0>0$, $(\Lambda_i^0)$, $(M_i)$.
\STATE $R\gets R_0$, $\delta\gets\rho/3$,
       $\Lambda_i\gets\max\{\delta,\Lambda_i^0\}$ for all $i$.
\WHILE{true}
  \STATE $D\gets C\cap\ball_R(s)$.
     Sort the pairs and compute $T_i,s_i,c_i$ using \eqref{eq:pa-parameters}.
  \STATE $U_i\gets s$ for all $i$.
     $y\gets\pass(1,1,0,\varnothing;U)$.
  \IF{the pass rejects pair $i$ with threshold $\tau_i$}
    \STATE $\Lambda_i\gets\max\{2\Lambda_i,\tau_i\}$ using the original pair index.
  \ELSE
    \IF{$\norm{y-s}\le R/2$}
      \RETURN $y$.
    \ENDIF
    \STATE $R\gets2R$.
  \ENDIF
\ENDWHILE
\STATE \textbf{procedure} $\pass(k,\beta,b,\mathcal M;U)$
\IF{$k=n+1$}
  \RETURN $\displaystyle\arg\min_{z\in D}\sum_{m\in\mathcal M}m(z)$ using the projection formula above.
\ENDIF
\STATE $y\gets U_k$.
\FOR{$t=0,\ldots,T_k-1$}
  \STATE $u\gets U_k$, $\alpha\gets\alpha_t$, $\beta'\gets\beta\alpha$,
     $b'\gets b+\beta(1-\alpha)y$.
  \STATE $x\gets\beta'u+b'$, $g\gets\nabla p_k(x)$,
     $q\gets Q_k(u)$, $H\gets\beta'c_k$.
  \STATE $m(z)\gets(H/2)\norm{z-u}^2+\ip{g+q}{z}$.
  \STATE $v\gets\pass(k+1,\beta',b',[\mathcal M,m];U)$.
     Propagate any rejection immediately.
  \STATE $x'\gets\beta'v+b'$.
     Compute $\tau_k$ using \eqref{eq:secant-test}.
  \IF{$\tau_k>\Lambda_k$}
    \RETURN a rejection for pair $k$ with threshold $\tau_k$.
  \ENDIF
  \STATE $y\gets(1-\alpha)y+\alpha v$,
    $U_k\gets\Pi_D\bigl(v+(q-Q_k(v))/H\bigr)$.
\ENDFOR
\RETURN $y$.
\end{algorithmic}
\end{algorithm}

Секущая в проверке вычисляется по формуле
\begin{equation}\label{eq:secant-test}
 \tau_k=\frac{2\ip{\nabla p_k(x')-\nabla p_k(x)}{x'-x}}{\norm{x'-x}^2},
 \qquad \tau_k=0\quad\text{при }x'=x.
\end{equation}

\subsection{Гарантия и цена адаптации}

Пусть $r_\star=\norm{s-z^\star}$,
$\bar R=\max\{R_0,24r_\star\}$ и
\[
 A_i^{\rm ball}(R)=
 \max_{\xi\in C\cap\ball_R(s)}(L_{0,i}+L_{1,i}\norm{\nabla p_i(\xi)}),
 \qquad
 \bar\Lambda_i=\max\{\lambda_i,4A_i^{\rm ball}(\bar R)\}.
\]
Эти величины используются в анализе. Алгоритму достаточно стартового
радиуса, желаемой точности, $M_i$ и доступа к проекциям.

\begin{theorem}[Путь А]\label{thm:patha}
Алгоритм завершается с $z_{\rm out}\in C$ и гарантией
\eqref{eq:relative}. Все попытки имеют радиус не больше $\bar R$
и кривизны $\Lambda_i\le\bar\Lambda_i$. Число попыток $N$, включая
успешную, удовлетворяет
\begin{equation}\label{eq:pa-restart}
 N\le 1+\left\lfloor\log_2(\bar R/R_0)\right\rfloor
       +\sum_i\left\lfloor\log_2(\bar\Lambda_i/\lambda_i)\right\rfloor .
\end{equation}
В исходной нумерации для обоих оракулов пары $i$ имеем
\begin{equation}\label{eq:pa-cost}
 \max\{N_{\nabla p_i},N_{Q_i}\}\le
 6\cdot17^{n-1}N
 \max\left\{1,\sqrt{3\Lambda_i^{\rm final}/\rho},\,3M_i/\rho\right\}.
\end{equation}
\end{theorem}

При $L_{1,i}=0$ можно взять $A_i^{\rm ball}(R)=L_{0,i}$.
С точностью до констант по глубине и множителя числа попыток
восстанавливается классическая зависимость по паре. При
$L_{1,i}>0$ наблюдаемая $\Lambda_i^{\rm final}$ может быть умеренной,
но верхняя оценка $A_i^{\rm ball}(\bar R)$ в общем случае может расти
экспоненциально с $L_{1,i}\bar R$. В формуле \eqref{eq:pa-cost}
оператор и градиент одной пары используют общий максимум.

\subsection{Почему алгоритм работает}

\subsubsection{Сертификат одного шага}

Выпуклость ограничивает остаток линеаризации $p_k$ сверху величиной
\[
 \ip{\nabla p_k(x')-\nabla p_k(x)}{x'-x}.
\]
Поэтому после проверки
$(\tau_k/2)\norm{x'-x}^2$ покрывает этот остаток. В координатах
уровня $f_k(z)=p_k(\beta'z+b')/\beta'$ достаточно кривизны
$L=\beta'\tau_k$. Из $\beta'\ge s_k$ следует $H\ge L+2M_k$.
Для сравнения с $w\in D$ пишем $u^+=\Pi_D(v+(Q_k(u)-Q_k(v))/H)$.
Выпуклость, монотонность и выделение квадрата дают
\begin{align}
 m_k(v)-m_k(w)+\frac H2\norm{v-w}^2
 &\ge f_k(v)-f_k(w)+\ip{Q_k(w)}{v-w}\notag\\
 &\quad+\frac H2\bigl(\norm{u^+-w}^2-\norm{u-w}^2\bigr).
 \label{eq:pa-model}
\end{align}
При выделении квадрата теряется не больше
$M_k^2\norm{v-u}^2/(2H)$. Этот член покрывает запас $H-L$.

\subsubsection{Сложение по уровням}

Зазор смешанной задачи на уровне $k$ имеет вид
\begin{align*}
 \mathcal G_k(z,w)
 &=\sum_{j<k}\left(m_j(z)-m_j(w)+\frac{H_j}{2}\norm{z-w}^2\right)\\
 &\quad+\sum_{i\ge k}\left(
 \frac{p_i(\beta z+b)-p_i(\beta w+b)}{\beta}
 +\ip{Q_i(w)}{z-w}\right).
\end{align*}
В листе он неположителен по условию оптимальности минимума суммы
моделей. На шаге $y^+=(1-\alpha)y+\alpha v$ из выпуклости
и \eqref{eq:pa-model} следует
\[
 \frac{\mathcal G_k(y^+,w)-(1-\alpha)\mathcal G_k(y,w)}{\alpha}
 +\frac H2\bigl(\norm{u^+-w}^2-\norm{u-w}^2\bigr)
 \le\mathcal G_{k+1}(v,w).
\]
Подстановка оценки дочернего уровня и деление ещё на $\alpha$
показывают, что каждому центру соответствует разность квадратов
расстояний. Её коэффициент $(\beta c_i/2)\prod_{j=k+1}^{i}a_j^2$
постоянен вдоль цикла.
Равенство $(1-\alpha_{t+1})/\alpha_{t+1}^2=1/\alpha_t^2$
сокращает соседние зазоры. Умножение на $a_k^2$ после суммирования
даёт оценку для всего вызова уровня $k$. В корне получаем
\begin{equation}\label{eq:pa-energy}
 G(y,w)\le \frac12\sum_i \ell_i
     \bigl(\norm{s-w}^2-\norm{U_i^{\rm out}-w}^2\bigr),
 \qquad \ell_i=s_i^2c_i.
\end{equation}
В частности, эта оценка действует на любом фиксированном выпуклом
рабочем множестве после принятого прохода без проверки выходного радиуса.

\subsubsection{Расписание и цена прохода}

Обратные веса растут за шаг на величину от $1/2$ до $1$,
поэтому $1/T\le\alpha_{T-1}\le2/T$.
Из округления $\prod_{j\le i}T_j\le 3\cdot17^{i-1}h_i$.
Если $a_1h_1=b_1$, то $1/3\le b_1\le1$ и
$s_i h_i\le b_1 8^{-(i-1)}$. Это даёт
$\ell_1\in[\delta b_1^2,3\delta]$ и
$\ell_i\le7\cdot8^{-(i-1)}\ell_1$ для $i>1$.
Значит, $\sum_{i>1}\ell_i\le\ell_1$ и $\frac12\sum_i\ell_i\le3\delta$.
Каждый шаг требует не больше двух вызовов каждого оракула своей пары.
Лист требует одной проекции, как и каждая операторная поправка.
Эти оценки действуют и для оборванной попытки. Суммирование
даёт \eqref{eq:pa-cost}.

\subsubsection{Конечность поиска и глобальный зазор}

Если $z^\star\in D$, то подстановка $w=z^\star$ в корневую оценку
для префикса даёт $\norm{U_1^t-z^\star}\le\sqrt2 r_\star$.
Это следует из неотрицательности зазора и
$\sum_{i>1}\ell_i\le\ell_1$. Из нерастяжимости проекции
$\norm{u^+-v}\le\norm{v-u}/2$, откуда
$\norm{v-z^\star}\le2\norm{u^+-z^\star}+\norm{u-z^\star}$.
Следовательно, выход проходит проверку половины радиуса при
$R\ge12r_\star$. Вместе с удвоением это даёт границу $\bar R$.
На любом пробном отрезке $\tau_k\le 2A_i^{\rm ball}(\bar R)$
для соответствующей исходной пары. Обновляемая координата растёт
хотя бы вдвое и остаётся под своим пределом. Отсюда
\eqref{eq:pa-restart}.

На рабочем шаре \eqref{eq:relative} следует из \eqref{eq:pa-energy}.
Для произвольного $w\in C$ с $d=\norm{w-s}>R$ возьмём
$w'=(1-\theta)y+\theta w$ с $\theta=R/(2d-R)$.
Поскольку $\norm{y-s}\le R/2$, имеем $w'\in D$.
Выпуклость $p$ и монотонность $Q$ дают
$\theta G(y,w)\le G(y,w')$. Тогда
$G(y,w)\le\rho R^2/\theta=\rho R(2d-R)\le\rho d^2$.
